% GENERATED FILE. Edit canonical lessons, then run npm run generate:books.
% canonical-source-hash: fnv1a64:f28a00daa068b090
% canonical-lessons: RU-C390-soblyudat, RU-C390-obyazatelnyy, RU-C390-poryadok, RU-C390-pravo, RU-C390-otvetstvennost

\chapter{Keeping the rules}
\label{ch:ru-keeping-the-rules}
\hlchaptermodality{390}

\begin{chapteropening}
\textbf{By the end of this chapter:} \emph{I can say that we must keep the rules, that something is compulsory, that there is order, and talk about rights and responsibility.}

Complete the last lesson of chapter 390: I can say that we must keep the rules, that something is compulsory, that there is order, and talk about rights and responsibility.
\end{chapteropening}

\section[\texorpdfstring{\ru{соблюдать} (soblyudát)}{соблюдать (soblyudát)}]{\ru{соблюдать} (soblyudát) — to observe, to keep (a rule)}

\label{lesson:RU-C390-soblyudat}



\begin{quote}
Before the new one: say the Russian for permission, then the Russian for to let.
\end{quote}

\subsection*{You'll want to know: \ru{соблюдать}}
\textbf{\ru{соблюдать}} (\emph{soblyudát}) — \textquotedblleft{}to observe, to keep (a rule)\textquotedblright{}.

\begin{quote}
\ru{Все} \ru{должны} \ru{соблюдать} \ru{правила}.
\end{quote}

\begin{quote}
\emph{Vse dolzhný soblyudát právila.}
\end{quote}

\begin{quote}
Everyone must observe the rules.
\end{quote}

Imperfective; perfective \emph{\ru{соблюсти}}, rare. \emph{\ru{Соблюдать} \ru{закон}, \ru{диету}, \ru{тишину}}.

\begin{cousinweb}
\textbf{\ru{соблюдать}} is built from \emph{\ru{со}-} and \emph{\ru{блюсти}} \textquotedblleft{}to guard, to keep\textquotedblright{}, an old Slavic verb also behind \emph{\ru{наблюдать}} \textquotedblleft{}to observe\textquotedblright{}.
\end{cousinweb}

\begin{grammarlens}[title={What you've built}]
One more word for telling, explaining and describing at B1.
\end{grammarlens}

\subsection*{Your turn}
\begin{itemize}
\raggedright
  \item \emph{Say it:} \emph{soblyudát}
  \item \emph{Say it:} \emph{soblyudát}, once more
  \item \emph{Say it:} a sentence of your own with \emph{soblyudát}
  \item \emph{Recall:} say the Russian for permission, then the Russian for to let, then say \emph{soblyudát} again
\end{itemize}

\begin{tcolorbox}[breakable,colback=teal!4,colframe=teal!35!black,title={Before you move on}]
What does \ru{соблюдать} mean? (\textquotedblleft{}To observe, to keep (a rule)\textquotedblright{}.) Say it once more.
\end{tcolorbox}
\section[\texorpdfstring{\ru{обязательный} (obyazátelnyy)}{обязательный (obyazátelnyy)}]{\ru{обязательный} (obyazátelnyy) — compulsory, obligatory}

\label{lesson:RU-C390-obyazatelnyy}



\begin{quote}
Before the new one: say the Russian for to let, then the Russian for to observe.
\end{quote}

\subsection*{You'll want to know: \ru{обязательный}}
\textbf{\ru{обязательный}} (\emph{obyazátelnyy}) — \textquotedblleft{}compulsory, obligatory\textquotedblright{}.

\begin{quote}
\ru{Этот} \ru{экзамен} \ru{обязательный} \ru{для} \ru{всех}.
\end{quote}

\begin{quote}
\emph{Étot ekzámen obyazátelnyy dlya vsekh.}
\end{quote}

\begin{quote}
This exam is compulsory for everyone.
\end{quote}

Of a person it means reliable: \emph{\ru{он} \ru{очень} \ru{обязательный}}.

\begin{cousinweb}
\textbf{\ru{обязательный}} is built on \emph{\ru{обязать}} \textquotedblleft{}to oblige\textquotedblright{}, from \emph{\ru{о}-} and \emph{\ru{вязать}} \textquotedblleft{}to tie\textquotedblright{}, Proto-Slavic \emph{*vęzati}: what you are bound to do.
\end{cousinweb}

\begin{grammarlens}[title={What you've built}]
One more word for telling, explaining and describing at B1.
\end{grammarlens}

\subsection*{Your turn}
\begin{itemize}
\raggedright
  \item \emph{Say it:} \emph{obyazátelnyy}
  \item \emph{Say it:} \emph{obyazátelnyy}, once more
  \item \emph{Say it:} a sentence of your own with \emph{obyazátelnyy}
  \item \emph{Recall:} say the Russian for to let, then the Russian for to observe, then say \emph{obyazátelnyy} again
\end{itemize}

\begin{tcolorbox}[breakable,colback=teal!4,colframe=teal!35!black,title={Before you move on}]
What does \ru{обязательный} mean? (\textquotedblleft{}Compulsory, obligatory\textquotedblright{}.) Say it once more.
\end{tcolorbox}
\section[\texorpdfstring{\ru{порядок} (poryádok)}{порядок (poryádok)}]{\ru{порядок} (poryádok) — order}

\label{lesson:RU-C390-poryadok}



\begin{quote}
Before the new one: say the Russian for to observe, then the Russian for compulsory.
\end{quote}

\subsection*{You'll want to know: \ru{порядок}}
\textbf{\ru{порядок}} (\emph{poryádok}) — \textquotedblleft{}order\textquotedblright{}.

\begin{quote}
\ru{В} \ru{его} \ru{комнате} \ru{всегда} \ru{порядок}.
\end{quote}

\begin{quote}
\emph{V yegó kómnate vsegdá poryádok.}
\end{quote}

\begin{quote}
His room is always tidy.
\end{quote}

Masculine; genitive \emph{\ru{порядка}}. \emph{\ru{Всё} \ru{в} \ru{порядке}} — everything is all right. \emph{\ru{По} \ru{порядку}} — in order.

\begin{cousinweb}
\textbf{\ru{порядок}} is built from \emph{\ru{по}-} and \emph{\ru{ряд}} \textquotedblleft{}a row, a line\textquotedblright{}, inherited from Proto-Slavic \emph{*rędŭ}: order is things set in rows.
\end{cousinweb}

\begin{grammarlens}[title={What you've built}]
One more word for telling, explaining and describing at B1.
\end{grammarlens}

\subsection*{Your turn}
\begin{itemize}
\raggedright
  \item \emph{Say it:} \emph{poryádok}
  \item \emph{Say it:} \emph{poryádok}, once more
  \item \emph{Say it:} a sentence of your own with \emph{poryádok}
  \item \emph{Recall:} say the Russian for to observe, then the Russian for compulsory, then say \emph{poryádok} again
\end{itemize}

\begin{tcolorbox}[breakable,colback=teal!4,colframe=teal!35!black,title={Before you move on}]
What does \ru{порядок} mean? (\textquotedblleft{}Order\textquotedblright{}.) Say it once more.
\end{tcolorbox}
\section[\texorpdfstring{\ru{право} (právo)}{право (právo)}]{\ru{право} (právo) — a right; law (as a subject)}

\label{lesson:RU-C390-pravo}



\begin{quote}
Before the new one: say the Russian for compulsory, then the Russian for order.
\end{quote}

\subsection*{You'll want to know: \ru{право}}
\textbf{\ru{право}} (\emph{právo}) — \textquotedblleft{}a right; law (as a subject)\textquotedblright{}.

\begin{quote}
\ru{Каждый} \ru{имеет} \ru{право} \ru{на} \ru{отдых}.
\end{quote}

\begin{quote}
\emph{Kázhdyy iméyet právo na ótdykh.}
\end{quote}

\begin{quote}
Everyone has the right to rest.
\end{quote}

Neuter. \emph{\ru{Иметь} \ru{право} \ru{на}} + accusative, or + infinitive. It is also law as a subject: \emph{\ru{изучать} \ru{право}}.

\begin{cousinweb}
\textbf{\ru{право}} is built on \emph{\ru{правый}} \textquotedblleft{}right, correct\textquotedblright{}, inherited from Proto-Slavic \emph{*pravŭ} \textquotedblleft{}right, correct\textquotedblright{} (EtymDB): a right is what is due by what is correct.
\end{cousinweb}

\begin{grammarlens}[title={What you've built}]
One more word for telling, explaining and describing at B1.
\end{grammarlens}

\subsection*{Your turn}
\begin{itemize}
\raggedright
  \item \emph{Say it:} \emph{právo}
  \item \emph{Say it:} \emph{právo}, once more
  \item \emph{Say it:} a sentence of your own with \emph{právo}
  \item \emph{Recall:} say the Russian for compulsory, then the Russian for order, then say \emph{právo} again
\end{itemize}

\begin{tcolorbox}[breakable,colback=teal!4,colframe=teal!35!black,title={Before you move on}]
What does \ru{право} mean? (\textquotedblleft{}A right; law (as a subject)\textquotedblright{}.) Say it once more.
\end{tcolorbox}
\section[\texorpdfstring{\ru{ответственность} (otvétstvennost)}{ответственность (otvétstvennost)}]{\ru{ответственность} (otvétstvennost) — responsibility}

\label{lesson:RU-C390-otvetstvennost}



\begin{quote}
Before the new one: say the Russian for order, then the Russian for a right.
\end{quote}

\subsection*{You'll want to know: \ru{ответственность}}
\textbf{\ru{ответственность}} (\emph{otvétstvennost}) — \textquotedblleft{}responsibility\textquotedblright{}.

\begin{quote}
\ru{Родители} \ru{несут} \ru{ответственность} \ru{за} \ru{детей}.
\end{quote}

\begin{quote}
\emph{Rodíteli nesút otvétstvennost za detéy.}
\end{quote}

\begin{quote}
Parents bear responsibility for their children.
\end{quote}

Feminine. \emph{\ru{Нести} \ru{ответственность} \ru{за}} + accusative — to bear responsibility for.

\begin{cousinweb}
\textbf{\ru{ответственность}} is built on \emph{\ru{ответственный}} \textquotedblleft{}responsible\textquotedblright{}, from \emph{\ru{ответ}} \textquotedblleft{}an answer\textquotedblright{}, from \emph{\ru{от}-} \textquotedblleft{}back\textquotedblright{} and Old Russian \emph{větŭ} \textquotedblleft{}counsel\textquotedblright{}: having to answer for something.
\end{cousinweb}

\begin{grammarlens}[title={What you've built}]
One more word for telling, explaining and describing at B1.
\end{grammarlens}

\subsection*{Your turn}
\begin{itemize}
\raggedright
  \item \emph{Say it:} \emph{otvétstvennost}
  \item \emph{Say it:} \emph{otvétstvennost}, once more
  \item \emph{Say it:} a sentence of your own with \emph{otvétstvennost}
  \item \emph{Recall:} say the Russian for order, then the Russian for a right, then say \emph{otvétstvennost} again
\end{itemize}

\begin{tcolorbox}[breakable,colback=teal!4,colframe=teal!35!black,title={Before you move on}]
What does \ru{ответственность} mean? (\textquotedblleft{}Responsibility\textquotedblright{}.) Say it once more.
\end{tcolorbox}
